In \seamlessstreaming,
the encoder first generates hidden representations from streaming input audio, which are then fed into a simultaneous text decoder.
The simultaneous text decoder operates a policy,
which decides whether to predict the next token or pause the prediction to consume more context.
In our case,
we adopted Efficient Monotonic Multihead Attention (\emma) \citep{ma_efficient_2023} as the simultaneous policy.
EMMA is a monotonic attention-based method \citep{raffel_online_2017, chiu_monotonic_2018, arivazhagan_monotonic_2019, ma_monotonic_2019},
that estimates the monotonic alignment during the training time in an unsupervised fashion.
In \emma, each attention head operates an individual simultaneous policy.
For simplicity, we only discuss the algorithm for one head, which can be easily extended it to multi-lyaer multihead attention.

\subsubsection{Numerical stable estimation}
The policy in the simultaneous text decoder is parameterized by a stepwise probability $p_{i,j}$,
which represents the probability of generating the next prediction $\tgttext_{i}$
given partial input $\srcspeech_{\leq j}$ and partial output $\tgttext_{\leq i-1}$.
$p_{i,j}$ is computed through the stepwise probability networks.
\begin{equation}
	p_{i, j} = \texttt{Sigmoid}\left(\frac{ \texttt{FFN}_s(s_{i-1})^T \texttt{FFN}_h(h_j) + b }{\tau} \right)\!,
 \label{eq:streaming.stepwise_p}
\end{equation}
where $\texttt{FFN}_s$ and $\texttt{FFN}_h$ are multi-layer feedforward networks serving as energy projections,
$s_{i-1}$ is $i-1$-th decoder state and $h_j$ is $j$-th encoder states.
$b$ is a learnable bias, initialized by a negative value, which makes it an easier optimization from the offline policy.
 $\tau$ is the temperature factor to encourage polarized output from the stepwise probability network. 

 
To train the stepwise probability network,
we estimated the probability of the alignment of partial output $\tgttext_{\leq i-1}$ with partial input $\srcspeech_{\leq j}$, or the event where there happens to be $j$ source input when the partial output length is $i-1$.
Denote this probability as $\alpha_{i,j}$.
Given the stepwise probability from \Cref{eq:streaming.stepwise_p}, $\alpha_{i,j}$ can be represented as:  
\begin{equation}
    \label{eq:streaming.alpha_def}
    \alpha_{i,j} = p_{i,j} \sum_{k=1}^j \alpha_{i-1, k} \prod_{l=k}^{j-1}(1 - p_{i,l}),
\end{equation}
which is also known as \textit{monotonic attention}.

The computation of \Cref{eq:streaming.alpha_def} is not trivial in training time.
In prior work on monotonic attention,
\Cref{eq:streaming.alpha_def} was estimated into a closed-form representation \citep{raffel_online_2017}, 
which computed $\alpha$ in parallel.
However, such an estimation is numerically unstable and biased.
\emma, however,
introduces a numerical stable estimation of monotonic attention duration training time.
This expression can be reformulated into a parallel version \citep{ma_efficient_2023} \footnote{See \Cref{app:streaming_operators} for the definition of the operators}:
\begin{equation}
\alpha_{i,:} = p_{i,:} \odot \alpha_{i-1, :} \texttt{triu}_0\left(\texttt{cumprod}_2(1 - \texttt{triu}_1\left( J_{|\srcspeech| \times 1} \texttt{roll}_1(p_{i,:}) \right))\right)\!.
	\label{eq:emma}
\end{equation}
Notably, this estimation process is of closed-form, with the benefit of numerical stability and unbiasedness (as it does not require a denominator within the equation in \citep{raffel_online_2017}).
A comprehensive derivation of this closed-form estimation is provided in \Cref{app:emma_estimation}.

Furthermore, during training, we adapted the infinite-lookback \citep{arivazhagan_monotonic_2019, ma_monotonic_2019} version of monotonic attention. Once the $\alpha$ is estimated, we then estimated the softmax weights $\beta$ in encoder-decoder attention as
\begin{equation}
    \label{eq:streaming:milk}
   \beta_{ij} = \sum_{k=j}^{|\srcspeech|} \left( \frac{\alpha_{ik} e_{ij}}{\sum_{l=1}^k  e_{il}} \right)\!,
\end{equation}
where $e_{ij}$ is the attention energy between $j$-th input and $i$-th output.
\Cref{eq:streaming:milk} can be also computed in parallel as
\begin{equation}
     \beta_{i:} = e_{i:} \odot  \texttt{flip}_2 \left( \texttt{cumsum} \left( \texttt{flip}_2 \left(\alpha_{i:} \odot  \frac{1}{ \texttt{cumprod} \left( e_{i:} \right)} \right) \right) \right)\!.
\end{equation}
Finally, the attention of each head used in the training can be expressed as
\begin{equation}
    \text{Attention}(Q, K, V) = \beta V. 
\end{equation}

\subsubsection{Policy regularization}
Because only the infinite lookback variant of monotonic attention is applied to \seamlessstreaming, it is necessary to add regularization loss functions in order to prevent the model from learning a trivial offline policy. As such, we applied two regularizations to the monotonic attention.

\begin{description}[leftmargin=0cm]
    \item[Latency] describes how much partial information is needed for the model to start translating.
Consistent with prior work \citep{arivazhagan_monotonic_2019, ma_monotonic_2019},
we used expected delays for latency regularization.
Denoting the expected delays for $i$-th target text as $\bar{d}^{\text{text}}_i$,
which is computed as
\begin{equation}
    \bar{d}^{\text{text}}_i = \texttt{E}[j | i] = \sum_{k=1}^{|\srcspeech|} k \alpha_{i,k}.
\end{equation}
Given a latency metric $\mathcal{C}$, the loss term is then computed as
\begin{equation}
	\mathcal{L}_{\text{latency}} = \mathcal{C}(\bar{d}^{\text{text}}_1, \ldots, \bar{d}^{\text{text}}_{|\tgttext|}).
\end{equation}
\item[Variance] of the alignment characterizes the certainty of an estimation.
\citet{arivazhagan_monotonic_2019} proposed a method to reduce uncertainty by introducing a Gaussian noise to the input of stepwise probability network in \Cref{eq:streaming.stepwise_p}.
However, empirical results show that the method is inefficient, especially when used in speech translation models.
Therefore, we propose an alternative regularization-based strategy based on variance estimation.
Denoting $\bar{v_i}$ as the expected variances of the monotonic alignment for target token $\tgttext_i$, $\bar{v_i}$ can be expressed as
\begin{equation}
	\bar{v}_i = \texttt{E}[(j - \texttt{E}[j | i])^2 | i] = \texttt{E}[j^2 | i] -  \texttt{E}[j | i]^2 = \sum_{k=1}^{|\srcspeech|} k^2 \alpha_{i,k} - \left(\sum_{k=1}^{|\srcspeech|} k \alpha_{i,k}\right)^{\!\!2}\!.
\end{equation}
We then introduced the alignment variance loss as follows:
\begin{equation}
	\mathcal{L}_{\text{variance}} = \frac{1}{|\tgttext|}\sum_{i=1}^{|\tgttext|} \bar{v}_i.
\end{equation}
\end{description}
Finally, we optimized the model with the following objective:
 \begin{equation}
 	\mathcal{L}(\theta) = -\log p(\tgttext | \srcspeech) + \lambda_{\text{latency}} \mathcal{L}_{\text{latency}} +  \lambda_{\text{variance}} \mathcal{L}_{\text{variance}},
 \end{equation}
 where $\lambda_{\text{latency}}$ and $\lambda_{\text{variance}}$ are the loss weights.





